\chapter*{Abstract}
\addcontentsline{toc}{chapter}{Abstract}

Climate science relies on large raster datasets that are rarely accessible to
non-specialist users. Existing tools either require GIS expertise or offer only
static visualisations with no ability to compare locations or time periods
interactively.

This report describes the design and implementation of \textbf{Climatica}, a
web application developed during a 400-hour internship at the Universidad de
Valladolid. Climatica provides interactive access to WorldClim historical
climate data through four pages: a single-city climate profile with a
Walter-Lieth climatogram, a side-by-side city comparison, a multi-period
climate change explorer, and a regional heatmap renderer. City search is
backed by a locally hosted Apache Solr index built from GeoNames data, and
all external API calls are proxied through Next.js server-side route handlers
with Redis caching.

The system was implemented using Next.js~16, React~19, TypeScript, Leaflet,
and Recharts, and is delivered in three locales (English, Spanish, and
Ukrainian). Key technical contributions include a client-side Walter-Lieth
diagram built from scratch with SVG \texttt{clipPath} masking, a raster
heatmap engine that parses WorldClim pixel IRIs to reconstruct geographic
tile bounds, and a trigram-based city search pipeline backed by a GeoNames
index of over 12~million populated places.

\vspace{0.5cm}
\noindent\textbf{Keywords:} \palabrasClave

\chapter*{Acknowledgements}
\addcontentsline{toc}{chapter}{Acknowledgements}

I would like to thank my supervisors, Guillermo Vega Gorgojo and Eduardo
Gómez Sánchez, for their guidance, patience, and consistently insightful
feedback throughout this internship. Guillermo's expertise in semantic web
technologies shaped the architecture of this project in fundamental ways,
and Eduardo's thorough and detailed review of the application was invaluable
in identifying issues and guiding improvements.

% TODO: add acknowledgement for Natalia (role, surname)

I am also grateful to the research group at the Universidad de Valladolid for
welcoming me and providing a stimulating environment in which to work.

Finally, I thank my family and friends for their support during my time abroad
in Valladolid.

\tableofcontents
\listoffigures
\listoftables